\begin{definition}[{\cite[Chapter I, 2.3]{knus_qhfr}}] \label{definition:even_varepsilon_hermitian_form_on_a_module_over_a_ring_with_involution}
    Let $A$ be a \CrefAndHyperrefIfExist{definition:involution_on_a_ring}{ring with involution}, let $\varepsilon \in Z(A)$ satisfy $\varepsilon \cdot \overline{\varepsilon} = 1$, and let $M$ be an \CrefAndHyperrefIfExist{definition:module_of_a_ring}{(right) $A$-module}. Write $\operatorname{Sesq}_\varepsilon(M)$ and $\operatorname{Sesq}^\varepsilon(M)$ as in \Cref{notation:spaces_of_varepsilon_hermitian_forms_and_even_varepsilon_hermitian_forms_on_module_over_a_ring_with_involution}.

    One calls elements of $\operatorname{Sesq}_\varepsilon(M)$ \hldef{even $\varepsilon$-Hermitian forms}; in other words, a \CrefAndHyperrefIfExist{definition:sesquilinear_form_on_a_module_over_a_ring_with_involution}{sesqeuilinear form} $b \in \operatorname{Sesq}_{A}(M)$ is an even $\varepsilon$-Hermitian form if and only if there exists some sesquilienar form $k \in \operatorname{Sesq}_A(M)$ such that $b = k + \varepsilon k^*$(\Cref{definition:adjoint_map_of_a_sesquilinear_form_on_a_module_over_a_ring_with_involution}). 

    In general, $\Sesq_\varepsilon(M) \subset \Sesq^\varepsilon(M)$. If $2$ is a unit in $A$, then $\Sesq_\varepsilon(M) = \Sesq^\varepsilon(M)$, since $b = \frac{1}{2} S_{\varepsilon}(b)$ for $b \in \Sesq^\varepsilon(M)$. 
\end{definition}